% ============================================================================
\clearpage
\section*{v4.9.0-r17 Source-Only: Entropy between the Frames, the Photon as Information Singularity, and a Scale Dimension}
\addcontentsline{toc}{section}{v4.9.0-r17 Source-Only: Entropy between the Frames, the Photon as Information Singularity, and a Scale Dimension}
\label{sec:v490-r17-entropy-scale}
% ============================================================================

\paragraph{Purpose.}
The main body works from the photon frame.  The author asked, on 2026-10-22, three things.  Must several structures coincide
for that frame to arise at all?  Should the photon be treated as an information singularity?  Would a new degree of freedom
help, one in which conversions between frames can be parked?  This layer answers the second and third questions in the state
disk of r04 and records a coincidence that bears on the first.  The script is \texttt{llm/k56\_entropie\_skala.py}.  No value
is compared with data.  Earlier layers are not rewritten.

\noindent\textbf{Source-only status.}
This v4.9.0-r17 layer is source-only.  It does not build a PDF, does not assign a new DOI,
and does not authorize a public release.  The v4.9.0 Version DOI
\href{https://doi.org/10.5281/zenodo.23534954}{10.5281/zenodo.23534954} and the Concept DOI
\href{https://doi.org/10.5281/zenodo.19210728}{10.5281/zenodo.19210728} are unchanged by it.

\subsection*{v4.9.0-r17 the frame factor in the state disk}
\label{subsec:v490-r17-frame-factor}

In the state disk of r04 a field has the Bloch vector \(b\) with \(|b|=1/\cosh2\eta_f\), where \(\eta_f\) is the rapidity of the
observer against the frame in which \(E\parallel B\).  The state matrix \(\rho=\tfrac12(1+b\cdot\sigma)\) has the eigenvalues
\(\tfrac12(1\pm|b|)\).

\paragraph{Theoremlet v4.9.0-r17.1.}
For an observer offset by \(\eta\):
\[
|b|=\frac1{\cosh2\eta},\qquad
p_\pm=\frac12\Bigl(1\pm\frac1{\cosh2\eta}\Bigr),\qquad
p_+=\frac1{1+\beta^2},\qquad p_-=\frac{\beta^2}{1+\beta^2},\qquad \frac{p_-}{p_+}=\beta^2,
\]
with \(\beta=\tanh\eta\).  At the offset of the source, \(\cosh2\eta=\gL\), the frame factor is the inverse length of the Bloch
vector, \(\gL=1/|b|\).  At the Hopf value of the source, \(\beta=1/\sqrt3\) and \(\gL=2\), the eigenvalues are \(3/4\) and \(1/4\).
At the measured value they are \(0.76818\) and \(0.23182\).

\paragraph{Proof sketch.}
\(\cosh2\eta=(1+\beta^2)/(1-\beta^2)\), hence \(\tfrac12(1+1/\cosh2\eta)=1/(1+\beta^2)\).  Checked numerically.  \(\square\)

\paragraph{Consequence.}
Three numbers of the earlier layers are one.  The C5 factor \(3/4\) (r10, r11) is the larger eigenvalue of the doublet state at
the Hopf offset.  The exported quarter is the smaller one.  The lattice value \(0.7683\) of r13 is the larger eigenvalue at the
measured offset.  The frame factor \(\gL\) is the inverse length of the same Bloch vector.  The eigenvalue ratio is \(\bsym^2\), the
quantity that Theorem~\ref{thm:dark-energy} calls \(\Omega_M\).  The eigenvalue itself, \(0.2318\), is not \(\Omega_M\).

\subsection*{v4.9.0-r17 entropy as order parameter, and the photon}
\label{subsec:v490-r17-entropy}

\paragraph{Theoremlet v4.9.0-r17.2.}
The entropy \(S=-p_+\log_2p_+-p_-\log_2p_-\) of the doublet state increases monotonically from \(0\) in the frame with
\(E\parallel B\) to \(1\) bit as \(\eta_f\to\infty\).  It is \(0.811\) bit at the Hopf offset and \(0.781\) bit at the measured
offset.  The photon, \(b=0\), is the maximally mixed state.  In the Bures metric of the disk,
\(\mathrm ds^2=\tfrac14\bigl(\mathrm dr^2/(1-r^2)+r^2\mathrm d\theta^2\bigr)\), the radial distance from the state with \(E\parallel B\)
to the state at offset \(\eta\) is \(\tfrac12\operatorname{gd}(2\eta)\).  The photon lies at the finite distance \(\pi/4\), and the remaining
distance is \(\tfrac12\Pi(2\eta)\), half the angle of parallelism of r07 at twice the rapidity.

\paragraph{Proof sketch.}
\(r=\operatorname{sech}2\eta\) gives \(\mathrm dr^2/(1-r^2)=4\operatorname{sech}^22\eta\,\mathrm d\eta^2\), hence \(\mathrm ds=\operatorname{sech}2\eta\,\mathrm d\eta\)
and \(\int_0^\eta\operatorname{sech}2x\,\mathrm dx=\tfrac12\operatorname{gd}(2\eta)\), with \(\operatorname{gd}(\infty)=\pi/2\) and
\(\Pi=\pi/2-\operatorname{gd}\).  Checked numerically.  \(\square\)

\paragraph{The photon as information singularity, made precise.}
The metric itself is regular at the photon.  What becomes singular is the direction: \(r^2\,\mathrm d\theta^2\to0\), so the angle
\(\theta\), which carries the distinction between matter and tachyon (\(b_1\)) and the side of r02 (\(b_2\)), is undefined at
\(b=0\).  This is the polar-coordinate form of the branch point of r19.2.  In this sense the photon is an information
singularity.  It is the state of maximal entropy, in which the information that distinguishes the frames is absent.  It is
reached by no finite boost (Theoremlet r19.1), only in the limit, and at a finite information distance.

\paragraph{On the author's first question.}
The photon is the point where the two eigenvalues coincide, \(p_+=p_-=\tfrac12\), and where matter and tachyon lose their
distinction.  In the state disk the special frame is therefore the frame in which the structures that separate the other
frames fall together.  This supports the author's conjecture.  It also shows that the photon frame is not a frame like the
others: it is the limit point and reference of the disk.  A finite factor \(\gL\) ``to the photon frame'', as in
\S\ref{sec:phase-constants}, corresponds to the state with \(|b|=1/\gL\), not to \(b=0\).

\subsection*{v4.9.0-r17 a scale dimension for the conversions}
\label{subsec:v490-r17-scale}

The rule \(X^{(\mathrm{matter})}=\gL^{k}X^{(\mathrm{photon})}\) is a scaling with weights \(k\).  Carrying the scale as a coordinate
\(\sigma\) of its own parks the conversion there, \(X(\sigma)=e^{k\sigma}X(0)\), with the photon frame at \(\sigma=0\).

\paragraph{Theoremlet v4.9.0-r17.3.}
\begin{enumerate}[leftmargin=2em]
\item[(a)] \(\sigma\) is a coordinate in which conversions compose as shifts only if the factor is a homomorphism of the
rapidity: \(\sigma=\eta\), \(X=e^{k\eta}X(0)\), the Doppler law of way (C) in r01.  The Lorentz factor is not one,
\(\cosh\eta_1\cosh\eta_2\neq\cosh(\eta_1+\eta_2)\).  Way (A), \(\cosh2\eta=e^{\sigma}\), holds only at the tribonacci point (r01.3).
\item[(b)] Combinations of total weight \(0\) are independent of \(\sigma\), for example the Compton frequency
\(m_ec^2/\hbar\) of r08.2.  Every measurable ratio of weight \(w\neq0\) moves as \(e^{w\sigma}\).
\end{enumerate}

\paragraph{Proof sketch.}
(a) \(e^{\eta_1}e^{\eta_2}=e^{\eta_1+\eta_2}\); the Lorentz case numerically \(1.432\) against \(1.864\).  (b) additivity of \(k\).  \(\square\)

\paragraph{Reading.}
The author's parking dimension has a definite form.  It is a line on which frames are points and conversions are shifts.
Its coordinate is the rapidity, and the conversion law it allows is the Doppler law.  It also locates the photon: on the
\(\sigma\)-line of the source the photon frame sits at \(\sigma=0\), while in the state disk it sits at \(\eta_f=\infty\).  The two
pictures differ because \S\ref{sec:phase-constants} measures offsets from the photon, whereas the disk measures them from the
rest frame of the field.  Way (C) follows from the request that the parked conversions compose.  The observation of r07
(on-shell \(\cot\theta_W\) close to \(\gL\), way A) points the other way.  It is not a test, and decision 5 of r15 stays with
the author.

\subsection*{v4.9.0-r17 register entries}
\label{subsec:v490-r17-register}

\begin{center}
\small
\begin{tabular}{@{}>{\raggedright\arraybackslash}p{0.34\textwidth}>{\raggedright\arraybackslash}p{0.58\textwidth}@{}}
\toprule
statement & status after v4.9.0-r17\\
\midrule
frame factor \(\gL\) & inverse length of the Bloch vector at the offset, \(\gL=1/|b|\) (r17.1)\\
C5 factor \(3/4\), lattice value \(0.768\) & larger eigenvalue of the doublet state at the Hopf, respectively measured, offset (r17.1)\\
photon & maximally mixed state; information singularity as undefined direction at finite Bures distance \(\pi/4\); reached only in the limit (r17.2)\\
photon frame of \S\ref{sec:phase-constants} & reference and limit point; finite \(\gL\) corresponds to \(|b|=1/\gL\), not \(b=0\)\\
scale dimension \(\sigma\) & composable only as rapidity, giving the Doppler law (way C); weight-0 combinations invariant (r17.3)\\
\bottomrule
\end{tabular}
\end{center}

\subsection*{v4.9.0-r17 ledger and guard update}
\label{subsec:v490-r17-ledger-guard}

\[
R_{\rm info}^{v490\text{-}r17}=(0,0,1)
\]
for (eigenvalue identities established, photon as information singularity made precise, choice of conversion way left to the
author).

The layer spends no guard:
\[
R_{\rm glob}^{(v490\text{-}r17,+)}=(0,1,1,1,1),
\qquad
\Gamma_{\rm guard}^{v490\text{-}r17}=(1,1,1,1,1),
\qquad
\beta_4=1/256\ \text{unspent}.
\]
It does not imply exact finite \(C5\) zero closure, exact all-readout terminality,
\(C6/T2\), a global photon-clock/RH theorem, or global all-readout closure.

